%% kapa: exam=1 round=33 year=2022 period=1 form=A subject=real_estate subject_ko=부동산학원론 seq=3 number=83
%% source: Q-Net 감정평가사 기출문제 (cst00309) artlSeq=5209997 file=raw/1008/5209997/2022년도 제33회 감정평가사 1차시험 1교시(민법,경제학,부동산학).pdf
%% conversion: kordoc 4.24.51 -> tools/gen1.py (auto); review=proofread-sonnet; flags=none
\begin{question}{83}
토지의 특성에 관한 설명이다. ( )에 들어갈 내용으로 옳게 연결된 것은?
\begin{framedtext}
○ ( ㄱ )은 토지에 대한 소유욕을 증대시키며 토지이용을 집약화시킨다. \par
○ ( ㄴ )은 임장활동과 지역분석의 근거가 된다. \par
○ ( ㄷ )은 토지간의 비교를 어렵게 하며 완전한 대체를 제약시킨다.
\end{framedtext}
\begin{choices}
  \item ㄱ: 개별성, ㄴ: 부동성, ㄷ: 영속성
  \item ㄱ: 영속성, ㄴ: 부동성, ㄷ: 용도의 다양성
  \item ㄱ: 영속성, ㄴ: 인접성, ㄷ: 용도의 다양성
  \item ㄱ: 부증성, ㄴ: 인접성, ㄷ: 부동성
  \item ㄱ: 부증성, ㄴ: 부동성, ㄷ: 개별성
\end{choices}
\end{question}
\answer{5}
